% TOP_PROBLEM: {"id": 2303002, "title": "Research Problems in Function Theory — Problem 3.2", "category_id": 8, "category": {"id": 8, "name": "analysis", "display_name": "Analysis", "description": "Limits, continuity, calculus, and function theory.", "slug": "analysis", "order_index": 8, "created_at": "2026-07-31T15:26:25.671Z"}, "status": "open"}
% FIRST_POSTED: null
% ENTRY_KIND: statement_only
% Imported from: batch12.tex; UnsolvedMath ID 2303002
% Compile twice: lualatex 2303002.tex
\documentclass[11pt]{article}
\usepackage{fontspec}
\setmainfont{Latin Modern Roman}
\usepackage{amsmath,amssymb,mathtools,unicode-math}
\setmathfont{Latin Modern Math}
\usepackage{xurl,hyperref}
\hypersetup{colorlinks=true,urlcolor=blue,pdftitle={UnsolvedMath Problem 2303002}}
\newcommand{\sref}[2]{\hyperlink{source:#1:#2}{[S#2]}}
\newcommand{\cataloguescope}[1]{\paragraph{Mathematical question.}#1}
\begin{document}
% UnsolvedMath ID 2303002; source catalogue code AMR-022-3002
\setcounter{section}{2303001}
\section{Research Problems in Function Theory — Problem 3.2}\label{2303002}
{\small\sffamily UnsolvedMath ID 2303002 · Analysis}
\subsection{Definitions and mathematical statement}

\paragraph{Shared notation.}$\mathbb N=\{1,2,\ldots\}$, and $\mathbb Z,\mathbb Q,\mathbb R,\mathbb C$ denote the integers, rationals, reals and complex numbers. Any other conventions are specified locally.


A subharmonic function is an upper semicontinuous function with values in $[-\infty,\infty)$, not identically $-\infty$ on a connected component, that satisfies the mean-value inequality on every closed ball in its domain. A real harmonic function is a twice continuously differentiable function $u$ satisfying $\Delta u=0$. A path to infinity in $\mathbb R^m$ is a continuous $\gamma:[0,\infty)\to\mathbb R^m$ with $|\gamma(t)|\to\infty$.
The source update attributes an affirmative answer to Fuglede and explains that smooth or piecewise-linear paths can be chosen.
\paragraph{Statement recorded in the supplied source.}
For every $m\geq3$ and every nonconstant harmonic function $u:\mathbb R^m\to\mathbb R$, is there a path $\gamma$ to infinity such that $u(\gamma(t))\to+\infty$? \sref{2303002}{2}
\subsection{Short English statement}
Does every nonconstant harmonic function in dimension at least three tend to positive infinity along a continuous path?
\subsection{Sources}
\begin{description}
\item[\hypertarget{source:2303002:1}{S1}] ULAM.AI and the UnsolvedMath Contributors, UnsolvedMath: A Curated Collection of Open Mathematics Problems, record 2303002 (AMR-022-3002). CC BY 4.0. \url{https://huggingface.co/datasets/ulamai/UnsolvedMath}.
\item[\hypertarget{source:2303002:2}{S2}] W. K. Hayman and E. F. Lingham, \emph{Research Problems in Function Theory}, arXiv:1809.07200v2 (2018), Problem 3.2 and Update 3.2. \url{https://arxiv.org/abs/1809.07200}.
\end{description}
\label{end:2303002}
\end{document}
